%%%%%%%%%%%%Outline%%%%%%%%%%%%%%%%%%

%%%
% message<xx>
% xx

%%%%%%%%%%%%Outline%%%%%%%%%%%%%%%%%%

\section{Conclusion}
\label{sec:conclusion}

We presented \textsc{BpfReJIT}, a dynamic, minimally invasive, and
extensible compilation framework that brings runtime-guided optimization
to kernel-resident eBPF programs. Our key insight is a
microkernel-inspired separation of concerns: the kernel defines
\emph{what can be optimized} through minimal extensions (${\sim}$600
lines of code), while a userspace daemon defines \emph{how to optimize}
through compilation passes informed by runtime profiling---analogous to
how JVM HotSpot and LLVM achieve high performance through extensible,
feedback-directed compilation.

\textsc{BpfReJIT} is fully transparent to existing eBPF tools, loaders,
and subsystems. The BPF verifier provides identical safety guarantees
for recompiled programs as for freshly loaded ones, ensuring that
kernel safety is never compromised regardless of daemon behavior. The
kinsn mechanism enables platform-specific instruction extensions via
kernel modules with zero verifier modifications, providing an
LLVM-pass-like extensibility model for the OS kernel.

\textit{[TODO: Add final evaluation summary numbers for v2 once
evaluation is complete.]}

\textsc{BpfReJIT} provides a practical, upstream-compatible path toward
dynamic, extensible compilation frameworks in operating system kernels.